\subsection{What is supplied and what was checked}

The archive contains this PDF and all its LaTeX inputs, the principal
integer polynomial and its 319-variable companion, both natural quartics,
the three compact positive circuits, and both eleven-integer circuits.
Every endpoint has its ordered variable list, literal coefficients or
complete arithmetic DAG, source binding, and an evaluator or exact
verification entry point. Imported source proofs and licenses retain
their provenance.

The top-level command
\begin{center}\vtt{python3 -B verify_all.py --rebuild}\end{center}
runs the independent component checks and deterministic rebuilds from the
archive alone, in ordinary and optimized Python. Python's standard library
is sufficient. No network, proof search, machine simulation, symbolic
algebra package, or original repository checkout is required. The PDF
build is separate and uses pdfLaTeX through the supplied build script.

The finite checks cover the following distinct obligations:
\begin{itemize}
\item The direct integer source changes are checked literally, including
the nine powers at base $q+1$, the remaining fourteen guards, and the two
new square-coordinate rows. The 217 monic elimination steps, translations,
and all 6,607 collected coefficients of the principal equation are replayed.
The parity companion has its own complete replay.
\item The natural branch checks every changed and unchanged raw row,
both simultaneous coordinate shifts, all 101 elimination certificates per
endpoint, and all quartic coefficients. Its isolated producer regenerates
the twelve normative JSON files exactly. Four corrupted artifacts per
endpoint are rejected.
\item The positive branch authenticates the entire fixed coefficient tuple,
checks ten exact free-polynomial identities, every graph reference and live
operation, global degree bounds, and complete univariate polynomial
specializations with nonzero leading coefficients. Twelve actual gate
mutations are rejected. The final component receipt records 276 checks.
\item The eleven-integer branch reconstructs all six fixed coefficients
from their finite recipe, checks local identities and every live gate,
proves the full leading monomial over the coefficient ring, and rejects
literal, input, and degree corruptions.
\end{itemize}

These checks are not numerical evidence for consistency. They certify
finite algebra, exact artifacts, and transformations. The unbounded
Pell, factorial, CRT, and proof-calculus arguments are the ordinary
mathematical proofs in this report and the identified inherited sources.
Separate agents reviewed the new sign-sensitive arguments; those reviews
are internal independent checks in this AI-assisted workflow, not outside
human endorsements. The work does not claim a newly compiled Lean/Rocq
development or a sealed object-language ZFC proof.

\subsection{Comparison and scope}

The repository was inspected at
\begin{center}\small
\vtt{1c5367452a0bd4dd7e9e4bceb1508b62dbd2ff61}.
\end{center}
It includes packages 127--133. This matters because an older summary table
still listed weaker endpoints. The comparisons in Section 1 use literal
equations and live circuits from the newest relevant packages, not the
stale table. A complete manuscript or equation from an earlier package is
credited wherever its proof interface is reused.

The direct 286-variable equation is the recommended unrestricted-integer
result of this iteration. It reduces dimension substantially while keeping
every coefficient below $2^{17}$ in absolute value. It is not the
fewest-variable quartic in the repository. The natural 235-variable form
has a different domain and much larger coefficients; it must not be
silently presented as an equation in 235 unrestricted integers. Likewise,
the 80-operation compact form has ten \emph{positive} inputs and larger
degree than the 83-operation form. The ordinary ten-integer constructions
using Sun's reduction remain a different dimension/degree/constant-size
tradeoff~\cite{sun}.

The source repository is unchanged by this work. The deliverable is an
integration-ready archive, with a suggested destination in the active
Diophantine consistency research collection and a source ledger that pins
all comparisons.

\subsection{Further research questions}

\begin{enumerate}
\item \textbf{Can other arithmetic guards be recovered semantically?}
The principal saving came from using a power certificate twice: once for
its value and once for its exponent's sign. A systematic dependency audit
of the remaining fourteen guards could find further removals. Each
candidate must be proved before the step that needs it; a circular
factorial or radix argument would introduce spurious roots.
\item \textbf{Can the natural one-row cap be transferred to the signed
compiler?} The natural bound uses a nonnegative quotient essentially.
A signed analogue would need to recover that quotient's sign or replace
the bound with another exact arithmetic condition. Merely reusing the
same equation over $\Z$ is unsound.
\item \textbf{Can abscissa orientation be made sparser?} The pooled
$(x_i-1)^2$ terms buy 27 deleted sign witnesses but account for the
principal support increase. A smaller subset of principal orientations,
followed by propagation along shared power relations, may retain most of
the dimensional saving with fewer collected monomials.
\item \textbf{Can the 80-operation construction retain a lower degree?}
The $q^3$ normalization removes three gates but raises the degrees of
the Pell data. A different quotient coordinate or a compatible lower
degree carrier would improve this tradeoff. The leading forms in
Section~\ref{sec:compact} identify the precise terms to target.
\item \textbf{Can the new arithmetic interface be formalized independently
of the whole compiler?} The negative-exponent lemma, even-frame anchor,
linear cofinal bound, and free auxiliary-coordinate transport are compact
targets for Lean or Rocq. Proving these first would make subsequent
compiler changes easier to audit. A full internal-ZFC derivation remains
a separate task.
\end{enumerate}

The most useful general lesson of this iteration is constructive: before
adding a sign witness, inspect whether an already required arithmetic
relation and its size bound force that sign. The new equations make this
principle precise in an encoding of finite set-theoretic refutations.
